% ------------------------------------------------------------
% Page 73
% ------------------------------------------------------------

\begin{center}
    {\Large\bfseries PASSAGE DE LA SOIE A LA LAINE}
\end{center}

\vspace{0.55cm}

Jouxtant la « carderie de soie », il y avait donc le foulon exploité par David Saumade

\vspace{0.25cm}

En 1859 le foulon est ainsi décrit : \textit{« Pas de bâtiment, si ce n’est un petit toit sous
lequel sont situées les 2 pots à fouler d’une valeur 400f \hspace{0.5cm} 10\% \hspace{0.5cm} 40 f\\
Valeur du cours d’eau \hspace{2.1cm} 600f \hspace{0.5cm} 5\% \hspace{0.5cm} 30 f\\
Le bâtiment grange \hspace{2.4cm} 200 f \hspace{0.5cm} 5\% \hspace{0.5cm} 10 f. »}

\vspace{0.3cm}

\textit{« ...le foulon est mis en mouvement par une 2° roue également verticale et qui
reçoit une chute particulière... dans les mois les plus secs de l’été les 2 roues
motrices du foulon et des métiers ne peuvent marcher simultanément mais cet
inconvénient est presque nul pour cette usine car elle ne travaille guère que dans
les mois d’automne et d’hiver, époque où l’eau afflue toujours en quantité
suffisante ».}

\vspace{0.35cm}

La carderie de bourres de soie a donc fait place à la filature de laine dans les
années 1859-1861. En 1859 l’exploitant de la « carderie » s’appelait \textbf{Lambert}
mais sur les 7 métiers, 3 avaient été démontés et venaient d’être remplacés par
des métiers à filer la laine et leurs accessoires.

\vspace{0.35cm}

En effet, c’est l’époque où s’est répandue la pébrine, maladie des vers à soie qui
provoqua la ruine de la soierie française et probablement la disparition de
l’élevage des vers à soie dans la région de Meyrueis, d’autant plus que la
fraîcheur des printemps demandait un chauffage onéreux des magnaneries.

\vspace{0.35cm}

\textbf{David Saumade}, exploitant du foulon, achètera le Moulin en 1861.

\vspace{0.35cm}

Il exploitera la filature avec plus ou moins de bonheur, au moins jusqu’à la mort
de sa femme, puisqu’en mars 1888 le moulin fut hypothéqué :

\vspace{0.25cm}

\begin{flushleft}
\hspace{0.35cm}
\parbox{0.88\textwidth}{
\textit{
« ... au profit de M. Jules Germain, filateur domicilié à Florac, en qualité de
légataire universel de Marie Eugénie Germain, décédée, femme du vendeur
pour garantie de reprises dotales de cette dernière se portant à la somme de
2.000 Fr. outre les intérêts conservés par la loi... »
}
}
\end{flushleft}

\vspace{0.25cm}

Incapable de régler cette somme \textbf{Saumade} dut vendre pour « purger cette
hypothèque ». (rembourser à son beau père la dot que ce dernier avait donné à sa
fille lors de son mariage !!!).
